% =============================================================================
% Chapter 17: Conclusion — ML Systems Lecture Slides (Volume II)
% =============================================================================
% This is the final chapter of Volume II. It synthesizes the six principles,
% reviews competencies mastered, and looks forward to compound AI systems.
% Target: 22 frames. Heavy on retrieval practice and reflection.
%
% IMAGES (SVGs converted to PDF):
%   - fleet-stack-synthesis.pdf    - six-principles.pdf
%   - principle-tensions.pdf       - competency-framework.pdf
%   - compound-ai-systems.pdf      - fermi-brain-vs-cluster.pdf
%   - post-silicon-frontier.pdf
% =============================================================================
\documentclass[aspectratio=169, 12pt]{beamer}
\usepackage{../../assets/beamerthememlsys}

\mlsyssetup{
  volume       = {Volume II},
  chapter      = {Chapter 17},
  logo         = {../../assets/img/logo-mlsysbook.png},
  instlogo     = {../../assets/img/logo-harvard.png},
  chaptertitle = {Conclusion},
}

% --- Fonts ---
\usepackage{fontspec}
\setsansfont{Helvetica Neue}[
  BoldFont={Helvetica Neue Bold},
  ItalicFont={Helvetica Neue Italic},
  BoldItalicFont={Helvetica Neue Bold Italic},
]
\setmonofont{JetBrains Mono}[Scale=0.85]

% --- Packages ---
\usepackage{booktabs}
\usepackage{amsmath}

% --- Image paths ---
\graphicspath{
  {images/}
}

% --- Helper: safe image include ---
\newcommand{\safeimg}[2][width=\textwidth,keepaspectratio]{%
  \IfFileExists{images/#2}{\includegraphics[#1]{#2}}{%
    \IfFileExists{#2}{\includegraphics[#1]{#2}}{%
      \fbox{\parbox[c][2.5cm][c]{0.85\linewidth}{\centering\footnotesize\textcolor{midgray}{[Missing image]}}}%
    }%
  }%
}

% --- Section count for navigation (must match actual \section{} count) ---
\setcounter{mlsystotalsections}{6}

\title{Conclusion}
\author{Vijay Janapa Reddi}
\institute{Harvard University}
\date{}

\begin{document}

% =============================================================================
% TITLE SLIDE
% =============================================================================
\mlsystitle{Conclusion}{Engineering Intelligence at Scale}{cover_conclusion.png}

% =============================================================================
% LEARNING OBJECTIVES
% =============================================================================
\begin{frame}{Learning Objectives}
\note{[2 min] Read through objectives. This is the synthesis chapter ---
everything connects. Ask: ``Which of the six principles do you remember
without looking at your notes?''}

\small
\begin{enumerate}
  \item Synthesize the \textbf{six principles} of distributed ML systems engineering
  \item Evaluate \textbf{interconnections} between infrastructure, training, fault tolerance, and operations
  \item Identify \textbf{principle tensions} and how engineers navigate trade-offs
  \item Apply the \textbf{Fleet Stack} as a professional engineering framework
  \item Assess the transition from the \textbf{Era of Scaling} to the \textbf{Era of Composition}
  \item Formulate strategies for building systems that serve humanity \textbf{responsibly}
\end{enumerate}

\end{frame}

% =============================================================================
\section{The Fleet Stack}
% =============================================================================

\begin{frame}{The Fleet Stack: Four Layers, One Discipline}
\note{[3 min] Walk through the four layers from bottom to top. Emphasize that
students have now built the complete stack. The arrow on the right captures
the sixth principle: scale creates qualitative change at every layer.
Ask: ``Which layer did you find most surprising?''}

% --- Layout: FULL-WIDTH diagram ---
\centering
\safeimg[width=0.88\textwidth,height=5.0cm,keepaspectratio]{fleet-stack-synthesis.pdf}

\vspace{0.1cm}
\mlsysinsight{Key:}{The complete stack equips engineers to make informed decisions at every level.}

\end{frame}

% --- ACTIVE LEARNING 1: Retrieval ---
\begin{frame}{Retrieval: Name the Six Principles}
\note{[2 min] Retrieval practice. Give students 60 seconds to write down all
six principles WITHOUT looking at notes. This tests their synthesis.
Walk around the room. Do NOT reveal yet --- next slide has the answer.}

\centering
\vspace{1.2cm}
{\Large\bfseries Close your notes.}

\vspace{0.6cm}
{\large Write down the \textbf{6 principles} of distributed ML systems.}

\vspace{0.6cm}
{\normalsize 60 seconds --- no peeking.}

\vspace{0.6cm}
{\small\textcolor{midgray}{Hint: think about the Fleet Stack layers.}}

\end{frame}

% =============================================================================
\section{Six Principles}
% =============================================================================

\begin{frame}{Six Principles of Distributed ML Systems}
\note{[3 min] Reveal after retrieval exercise. Walk through each principle
briefly --- students should recognize all of these from previous chapters.
Spend most time on principle 6 (scale creates change) since it is the
integrating principle. Ask: ``Which principle was hardest to remember?''}

% --- Layout: FULL-WIDTH diagram ---
\centering
\safeimg[width=0.92\textwidth,height=4.8cm,keepaspectratio]{six-principles.pdf}

\vspace{0.1cm}
\mlsysconcept{Pattern:}{Infrastructure constrains, operations expose, governance directs, scale transforms.}

\end{frame}

\begin{frame}{Principle Summary Table}
\note{[2 min] Reference table. Students can photograph this slide.
Each principle maps to a core question, a metric, and a chapter.
If short on time, skip this and point students to the textbook table.}

\footnotesize
\renewcommand{\arraystretch}{1.2}
\begin{tabular}{@{}lllp{3.2cm}@{}}
  \toprule
  \textbf{Principle} & \textbf{Core Question} & \textbf{Key Metric} & \textbf{Chapter} \\
  \midrule
  \textbf{1. Communication} & What is the bottleneck? & BW utilization & Collective Comm. \\
  \textbf{2. Failure} & How do we recover? & MTBF, ckpt overhead & Fault Tolerance \\
  \textbf{3. Infrastructure} & What is possible? & FLOPS, mem BW & Compute Infra. \\
  \textbf{4. Responsible} & Who is affected? & Fairness, audits & Responsible AI \\
  \textbf{5. Sustainability} & What is the cost? & kWh, carbon & Sustainable AI \\
  \textbf{6. Scale} & What breaks at 1000$\times$? & Scaling efficiency & Distributed Training \\
  \bottomrule
\end{tabular}

\vspace{0.15cm}
\mlsysalert{Remember:}{All six must be satisfied simultaneously --- no principle operates in isolation.}

\end{frame}

% =============================================================================
\section{Tensions}
% =============================================================================

\begin{frame}{Principle Tensions: The Art of Trade-offs}
\note{[3 min] The key insight: these principles sometimes CONFLICT.
Walk through each tension pair. Communication optimization increases failure
exposure. Sustainability limits raw performance. Responsible AI adds latency.
The art is finding designs that balance all six within acceptable trade-offs.}

% --- Layout: FULL-WIDTH diagram ---
\centering
\safeimg[width=0.92\textwidth,height=5.0cm,keepaspectratio]{principle-tensions.pdf}

\vspace{0.05cm}
{\small Navigating these tensions defines the \textbf{art} of distributed ML engineering.}

\end{frame}

% --- ACTIVE LEARNING 2: Discussion ---
\begin{frame}{Discussion: Which Trade-off is Hardest?}
\note{[3 min] Turn-and-talk. Students discuss in pairs for 90 seconds.
Cold-call 2--3 pairs. There is no single right answer --- the point is
that trade-offs depend on context (startup vs.\ hyperscaler, training vs.\
serving, regulated vs.\ unregulated domain).}

\centering
\vspace{0.8cm}
{\large\bfseries Turn and Talk \textcolor{midgray}{(90 seconds)}}

\vspace{0.6cm}
{\large A hyperscaler must choose between:\\[0.2cm]
\textcolor{computestroke}{\textbf{(A)}} Maximizing training throughput\\[0.1cm]
\textcolor{datastroke}{\textbf{(B)}} Minimizing carbon footprint\\[0.1cm]
\textcolor{errorstroke}{\textbf{(C)}} Ensuring fairness across user groups}

\vspace{0.5cm}
\alert{Which trade-off is hardest to navigate, and why?}

\end{frame}

\begin{frame}{No Principle Exists in Isolation}
\note{[2 min] Production systems must balance all six principles simultaneously.
The textbook chapters describe an integrated whole --- not independent modules.
Use the concrete example: communication optimization may require sync patterns
that increase failure exposure. This is the real engineering challenge.}

\footnotesize
\begin{columns}[T]
  \begin{column}{0.55\textwidth}
    In production, principles create \textbf{cross-cutting constraints}:

    \vspace{0.1cm}
    \begin{itemize}\setlength\itemsep{0pt}
      \item Communication optimization $\to$ more sync $\to$ failure exposure
      \item Sustainability $\to$ limits infrastructure choices
      \item Responsible AI $\to$ adds latency to serving
      \item Scale amplifies \emph{every} tension
    \end{itemize}

    \vspace{0.1cm}
    \begin{mlsyscard}{crimson}
    {\scriptsize The engineer must find designs that balance all six principles within acceptable trade-offs.}
    \end{mlsyscard}
  \end{column}
  \begin{column}{0.42\textwidth}
    \scriptsize
    \renewcommand{\arraystretch}{1.1}
    \begin{tabular}{@{}lc@{}}
      \toprule
      \textbf{Conflict} & \textbf{Resolution} \\
      \midrule
      Comm $\leftrightarrow$ Fault & Async ckpt \\
      Perf $\leftrightarrow$ Green & Carbon sched. \\
      Ethics $\leftrightarrow$ Latency & Offline audit \\
      \bottomrule
    \end{tabular}
  \end{column}
\end{columns}

\end{frame}

% =============================================================================
\section{Competencies}
% =============================================================================

\begin{frame}{Three Core Competencies Mastered}
\note{[3 min] Walk through the three competency pillars. Students should feel
confident they have developed skills in all three areas. Distributed systems
is the technical core; production ops is the practical skill; governance
distinguishes a systems engineer from a model developer.
Ask: ``Which competency do you feel strongest in?''}

% --- Layout: FULL-WIDTH diagram ---
\centering
\safeimg[width=0.92\textwidth,height=5.0cm,keepaspectratio]{competency-framework.pdf}

\vspace{0.05cm}
{\small These competencies span \textbf{all four layers} of the Fleet Stack.}

\end{frame}

% --- ACTIVE LEARNING 3: Exercise ---
\begin{frame}{Exercise: Map the Competency}
\note{[3 min] Give students 90 seconds. This is a synthesis exercise that
forces them to connect specific technical skills to the three competencies.
Answer: (1) Distributed Systems --- Ring AllReduce is a communication pattern.
(2) Production Operations --- continuous batching is a serving technique.
(3) Governance --- differential privacy is an ethics/privacy mechanism.}

\footnotesize
{\small\bfseries Match each skill to its competency:}

\vspace{0.15cm}
\begin{columns}[T]
  \begin{column}{0.55\textwidth}
    \textbf{Skills:}
    \begin{enumerate}\setlength\itemsep{0pt}
      \item Ring AllReduce bandwidth analysis
      \item Continuous batching optimization
      \item Differential privacy implementation
    \end{enumerate}

    \vspace{0.1cm}
    \textbf{Competencies:}
    \begin{itemize}\setlength\itemsep{0pt}
      \item[\textcolor{computestroke}{A}] Distributed Systems
      \item[\textcolor{routingstroke}{B}] Production Operations
      \item[\textcolor{datastroke}{C}] Governance \& Ethics
    \end{itemize}

    {\scriptsize\textcolor{midgray}{(90 seconds --- then compare with a neighbor)}}
  \end{column}
  \begin{column}{0.40\textwidth}
    \pause
    \begin{mlsyscard}{datastroke}
    \textbf{Solution:}\\[0.05cm]
    {\scriptsize
    1 $\to$ \textcolor{computestroke}{\textbf{A}} Distributed Systems\\
    2 $\to$ \textcolor{routingstroke}{\textbf{B}} Production Operations\\
    3 $\to$ \textcolor{datastroke}{\textbf{C}} Governance \& Ethics\\[0.05cm]
    Each maps to a Fleet Stack layer.
    }
    \end{mlsyscard}
  \end{column}
\end{columns}

\end{frame}

% =============================================================================
\section{The Path Forward}
% =============================================================================

\mlsysfocus{The Compound Capability Law}{%
$\text{Capability} \propto \text{Model}_{IQ} \times (\text{Tools} + \text{Context} + \text{Planning})^N$%
}

\begin{frame}{From Scaling to Composition}
\note{[3 min] This is the forward-looking content. The Era of Scaling is
flattening --- both Moore's Law and neural scaling laws are hitting power
limits. The next 100x must come from orchestration (10x), not hardware (4x)
or algorithms (2.5x). The unit of compute shifts from FLOP to Reasoning Chain.
Ask: ``What does it mean for your career if orchestration is the bottleneck?''}

% --- Layout: FULL-WIDTH diagram ---
\centering
\safeimg[width=0.92\textwidth,height=4.8cm,keepaspectratio]{compound-ai-systems.pdf}

\vspace{0.05cm}
\mlsysconcept{Insight:}{The next breakthrough is a smarter \textbf{System Architecture}, not a larger GPU.}

\end{frame}

\begin{frame}{The 100$\times$ Efficiency Decomposition}
\note{[2 min] Walk through the math. Hardware gives 4x (B200 vs H100).
Algorithms give 2.5x (sparsity + distillation). The remaining 10x MUST come
from orchestration --- compound AI systems, reasoning loops, dynamic retrieval.
This is why the skills taught in this textbook are the highest-leverage
investment for the next decade of AI.}

\footnotesize
\textbf{Where does the next 100$\times$ improvement come from?}

\vspace{0.1cm}
$$100\times = \underbrace{4\times}_{\text{Hardware}} \times \underbrace{2.5\times}_{\text{Algorithm}} \times \underbrace{10\times}_{\text{Orchestration}}$$

\vspace{0.05cm}
\scriptsize
\renewcommand{\arraystretch}{1.1}
\begin{tabular}{@{}llll@{}}
  \toprule
  \textbf{Source} & \textbf{Gain} & \textbf{Mechanism} & \textbf{Status} \\
  \midrule
  \textcolor{computestroke}{\textbf{Hardware}} & 4$\times$ & Next-gen silicon (B200 vs H100) & Diminishing returns \\
  \textcolor{routingstroke}{\textbf{Algorithm}} & 2.5$\times$ & Structured sparsity + distillation & Incremental \\
  \textcolor{datastroke}{\textbf{Orchestration}} & \textbf{10$\times$} & Compound AI systems & \textcolor{datastroke}{\textbf{Largest opportunity}} \\
  \bottomrule
\end{tabular}

\vspace{0.1cm}
\begin{mlsyscard}{crimson}
{\scriptsize The bulk of future scaling must come from \textbf{System Orchestration}: reasoning loops, tool-use, and dynamic retrieval delivering 10$\times$ more utility from the same FLOPs.}
\end{mlsyscard}

\end{frame}

\begin{frame}{The Fermi Estimate of Intelligence}
\note{[3 min] This is the capstone quantitative exercise. Walk through the
numbers. The machine is 2500x faster at raw ops but the brain is 350,000x
more energy efficient. The punchline: we have conquered scale. The challenge
for the next generation is efficiency. This connects directly to the post-silicon
frontier and sustainability principle.}

% --- Layout: FULL-WIDTH diagram ---
\centering
\safeimg[width=0.88\textwidth,height=4.5cm,keepaspectratio]{fermi-brain-vs-cluster.pdf}

\vspace{0.1cm}
\mlsysalert{Punchline:}{We have conquered \textbf{Scale}. The challenge is \textbf{Efficiency}.}

\end{frame}

\begin{frame}{Post-Silicon Frontiers}
\note{[2 min] The energy wall is the binding constraint. Moving from 10 pJ/bit
(electrical) to 0.01 pJ/bit (optical) is a 1000x efficiency gain. This is
required to scale to 1 million GPUs without consuming the global power grid.
The enduring principles from this textbook (bandwidth, fault tolerance,
governance) apply to any substrate. If short: summarize in 60 seconds.}

% --- Layout: FULL-WIDTH diagram ---
\centering
\safeimg[width=0.92\textwidth,height=5.0cm,keepaspectratio]{post-silicon-frontier.pdf}

\vspace{0.05cm}
{\small The principles from this textbook \textbf{endure} across substrates.}

\end{frame}

% =============================================================================
\section{Wrap-Up}
% =============================================================================

% --- ACTIVE LEARNING 4: Retrieval Practice ---
\begin{frame}{What Were the Key Ideas?}
\note{[2 min] Final retrieval practice. Students write for 90 seconds, no notes.
This is the most important retrieval moment of the entire course --- they are
synthesizing everything from Volume II. Walk around the room.}

\centering
\vspace{1.5cm}
{\Large\bfseries Close your notes.}

\vspace{0.8cm}
{\large Write down the \textbf{4 most important ideas} from this entire volume.}

\vspace{0.8cm}
{\normalsize\textcolor{midgray}{90 seconds --- this is your synthesis of the course.}}

\end{frame}

\begin{frame}{Key Takeaways}
\note{[2 min] Reveal. Walk through each bullet. These are the four takeaways
from the textbook's conclusion chapter. Emphasize that the transition from
single-machine to distributed is qualitative, not quantitative. The ethical
weight of engineering decisions is the final and most important point.}

\footnotesize
\begin{itemize}\setlength\itemsep{3pt}
  \item \textbf{Six principles} define distributed ML systems: communication dominance, routine failure, infrastructure determination, responsible engineering, sustainability constraints, and qualitative scale effects.

  \item \textbf{Single-machine to distributed is qualitative}: new phenomena emerge at scale (stragglers, network partitions, heterogeneity) requiring distinct engineering approaches.

  \item \textbf{Production ML is an integrated whole}: infrastructure, training, fault tolerance, operations, security, and governance are interconnected --- no component operates in isolation.

  \item \textbf{Engineering decisions carry ethical weight}: the systems we build affect billions of lives through recommendations, medical AI, climate models, and accessibility tools.
\end{itemize}

\end{frame}

\begin{frame}{References}
\note{[1 min] Point students to canonical papers for further study.}

\small
\mlsysref{Dean+12}{Dean \& Ghemawat. ``Large Scale Distributed Deep Networks.'' NeurIPS 2012.}
\mlsysref{Dubey+24}{Dubey et al. ``The Llama 3 Herd of Models.'' Meta AI, 2024.}
\mlsysref{Sergeev+18}{Sergeev \& Del Balso. ``Horovod: Fast and Easy Distributed Deep Learning.'' 2018.}
\mlsysref{Strubell+19}{Strubell et al. ``Energy and Policy Considerations for Deep Learning.'' ACL 2019.}
\mlsysref{Patterson+21}{Patterson et al. ``Carbon Emissions and Large Neural Networks.'' 2021.}
\mlsysref{Amodei+16}{Amodei et al. ``Concrete Problems in AI Safety.'' 2016.}

\end{frame}

% =============================================================================
% CLOSING REFLECTION (replaces "Next Lecture" for final chapter)
% =============================================================================

\begin{frame}{}
\note{[2 min] This is the closing moment of the entire course. Read the three
imperatives slowly. Let them land. Then show the attribution. Thank the
students for their work. This is not a lecture slide --- it is a send-off.}

\centering
\vspace{1.5cm}

{\Large\bfseries Build systems that scale.}

\vspace{0.4cm}

{\Large\bfseries Build systems that endure.}

\vspace{0.4cm}

{\Large\bfseries Build systems that serve humanity well.}

\vspace{1.0cm}

{\normalsize\textcolor{midgray}{\textit{Prof.\ Vijay Janapa Reddi, Harvard University}}}

\end{frame}

\end{document}
